% Bagean: metodhe
% Bab: bukti
% Pérangan: pola inferensi

\documentclass[../../../include/open-logic-section]{subfiles}

\begin{document}

\olfileid{mth}{prf}{pat}

\olsection{Pola-Pola Inferensi}

Bukti kasusun saka inferensi siji-siji. Nalika nggawe inferensi,
kita lumrahé nandhani kanthi tembung kaya ``mula'', ``dadi'', utawa
``kanthi mangkono''. Inferensi asring gumantung marang siji utawa
loro kasunyatan kang wis kasedhiya ing bukti: bisa anggepan, bisa uga
kesimpulan saka inferensi sadurungé. Supaya cetha, kita bisa mènèhi
tandha marang pratelan-pratelan mau lan nyebut endi kang digunakaké
ing inferensi anyar. Inferensi uga asring nerangaké \emph{ngapa}
kesimpulan anyar katut saka bab-bab sadurungé. Ana sawatara pola
inferensi kang kerep banget digunakaké ing bukti; kita bakal
ngrembug sawatara ing ngisor iki. Pola liya, kaya bukti kanthi
induksi, luwih rumit lan bakal dirembug mengko.

Kita wis ngrembug siji pola inferensi: njlentrehaké lan nrapaké
dhéfinisi. Nalika dhéfinisi dijlentrehaké, pratelan kang ngandhut
definiendum dipratelakaké maneh nganggo definiens. Upamané, anggepen
ing bukti wis ditemokaké $D = E$ (a). Banjur dhéfinisi $=$ kanggo
himpunan bisa ditrapaké kanggo nyimpulaké: ``Mula, miturut
dhéfinisi saka (a), saben !!{element} saka~$D$ uga !!a{element}
saka~$E$, lan kosok baliné.''

Kadhangkala mbingungaké: pambenaran inferensi ora tansah ditulis
nalika inferensi kuwi digawé, nanging bisa luwih dhisik. Anggepen
kita durung mbuktekaké $D = E$, nanging iki tujuan kita. Yèn $D = E$
kesimpulan kang dikarepaké, tujuwan iki bisa dipratelakaké maneh
kanthi dhéfinisi: kanggo mbuktekaké $D = E$, kita kudu mbuktekaké
yèn saben !!{element} saka~$D$ uga !!a{element} saka~$E$, lan
kosok baliné. Dadi bukti bisa awujud: (a) buktèkna saben
!!{element} saka~$D$ iku !!a{element} saka~$E$; (b) saben
!!{element} saka~$E$ iku !!a{element} saka~$D$; (c) mula saka
(a) lan (b) miturut dhéfinisi $=$, $D = E$. Nanging biasané ora
ditulis kaya mangkono. Kita bisa nulis:
\begin{quote}
Kita arep nuduhaké $D = E$. Miturut dhéfinisi~$=$, iki padha karo
nuduhaké yèn saben !!{element} saka~$D$ iku !!a{element}
saka~$E$, lan kosok baliné.

(a) \dots (bukti yèn saben !!{element} saka~$D$ iku !!a{element}
saka~$E$) \dots

(b) \dots (bukti yèn saben !!{element}
saka~$E$ iku !!a{element} saka~$D$) \dots
\end{quote}

\subsection{Migunakaké Konjungsi}

Mbokmenawa pola inferensi kang paling prasaja yaiku njupuk siji
pérangan saka konjungsi minangka kesimpulan. Tegesé, yèn kita
nganggep utawa wis mbuktekaké $p$ lan~$q$, kita kena nyimpulaké
~$p$ (uga~$q$). Inferensi dhasar iki asring ora perlu disebut.
Upamané, sawisé njlentrehaké dhéfinisi $D = E$, kita ngerti yèn
saben !!{element} saka~$D$ iku !!a{element} saka~$E$ lan kosok
baliné. Saka kéné kita bisa nyimpulaké yèn saben !!{element}
saka~$E$ iku !!a{element} saka~$D$; iki pérangan ``kosok baliné''.

\subsection{Mbuktekaké Konjungsi}

Kadhangkala kang dijaluk dibuktèkaké awujud konjungsi:
``buktèkna $p$ lan $q$''. Ing kéné cukup nindakaké rong bab:
mbuktekaké $p$, banjur mbuktekaké $q$. Bukti bisa dipérang dadi
rong pérangan lan diwènèhi tandha supaya cetha. Nalika nyusun
cathetan wiwitan, panjenengan bisa nulis ``(1) Buktèkna $p$''
ing ndhuwur kertas lan ``(2) Buktèkna $q$'' ing tengah kertas.
Kadhangkala pitakonan ora kanthi gamblang njaluk bukti konjungsi,
nanging buktiné mbutuhaké iku. Upamané, kanggo mbuktekaké $D = E$,
sawisé njlentrehaké dhéfinisi~$=$, panjenengan kudu mbuktekaké:
saben !!{element} saka~$D$ iku !!a{element} saka~$E$
\emph{lan} saben !!{element} saka~$E$ iku !!a{element} saka~$D$.

\subsection{Mbuktekaké Disjungsi}

Yèn kang dibuktèkaké awujud disjungsi, yaiku pratelan ``$p$
utawa $q$'', cukup nuduhaké yèn salah siji pérangané bener.
Nanging arang banget ana pérangan kang langsung katut saka
anggepan téoréma. Luwih kerep, anggepané dhéwé awujud disjungsi,
utawa kita nuduhaké yèn kabèh obyèk saka sawijining jinis nduwèni
salah siji saka rong sipat, nanging obyèk kang béda nduwèni
sipat kang béda. Ing kéné bukti mawa kasus migunani (delengen
ing ngisor iki).

\subsection{Bukti Kondisional}

Akèh téoréma awujud kondisional: yèn $p$ bener, mula $q$ uga
bener. Nata buktiné gampang: anggepen pérangan wiwitan kondisional,
ing kéné $p$, banjur buktèkna kesimpulan~$q$ saka anggepan iku.
Dadi yèn téoréma nyatakaké ``Yèn $p$ mula $q$'', bukti diwiwiti
kanthi ``anggepen $p$'' lan dipungkasi kanthi mbuktekaké~$q$.

Kondisional bisa dipratelakaké kanthi cara béda-béda. Saliyané
``Yèn $p$ mula $q$'', téoréma bisa nyatakaké ``$p$ mung yèn
$q$'', ``$q$ yèn $p$'', utawa ``$q$, angger $p$''.
Kabèh padha tegesé lan mbutuhaké anggepan $p$ banjur bukti~$q$.
Élinga yèn bikondisional ``$p$ yèn lan mung yèn (iff) $q$''
iku gabungan rong kondisional: yèn $p$ mula $q$, lan yèn $q$
mula~$p$. Mula cukup gawe rong bukti kondisional. Nanging
kadhangkala pratelan ``iff'' bisa dibuktèkaké kanthi ngrangkè
sawatara pratelan ``iff'', diwiwiti saka ``$p$'' lan dipungkasi
ing ``$q$''; ing kéné saben langkah kudu pancèn ekuivalensi.

\subsection{Pratelan Universal}

Migunakaké pratelan universal iku prasaja: yèn sawijining bab
bener kanggo kabèh obyèk, bab iku uga bener kanggo saben obyèk
tartamtu. Upamané, yèn hipotesis bukti yaiku $A \subseteq B$,
tegesé, sawisé dhéfinisi~$\subseteq$ dijlentrehaké, kanggo
saben $x \in A$ uga $x \in B$. Mula yèn wis ngerti $z \in A$,
kita bisa nyimpulaké $z \in B$.

Mbuktekaké pratelan universal bisa katon rada angel. Lumrahé
pratelan iku awujud: ``Yèn $x$ nduwèni sipat~$P$, mula uga
nduwèni sipat~$Q$'' utawa ``Kabèh $P$s iku $Q$s''.
Mesthi, wujudé ora tansah pas kaya iki, lan perlu latihan
kanggo ngerti apa persis kang kudu dibuktèkaké. Nanging asring
kita kudu mbuktekaké yèn saben obyèk mawa siji sipat uga
nduwèni sipat liya.

Kanggo mbuktekaké pratelan universal, kita ngenalaké jeneng
utawa variabel kanggo obyèk mawa sipat kapisan, banjur
nuduhaké yèn obyèk mau uga nduwèni sipat kapindho.
Tegesé, kanggo mbuktekaké bab kanggo \emph{kabèh}~$P$s,
kita mbuktekaké kanggo sawijining~$P$ kang \emph{sembarang}.
Jeneng kang dienalaké iku jeneng kanggo~$P$ kang sembarang.
Lumrahé jeneng iki aksara siji, manut pakulinan: $n$
kanggo bilangan asli, $!A$ kanggo !!{formula}s, $A$
kanggo himpunan, $f$ kanggo fungsi, lan sapanunggalané.

Kunciné yaiku njaga sipat umum ing sadawaning bukti. Wiwitané,
anggepen obyèk sembarang ``$x$'' nduwèni sipat~$P$, banjur
nuduhaké, mung saka dhéfinisi lan anggepan kang diidini, yèn
$x$ nduwèni sipat~$Q$. Amarga ora ana katrangan khusus babagan
$x$ saliyané sipat $P$, kita bisa nyimpulaké yèn saben obyèk
mawa $P$ uga nduwèni sipat~$Q$. Cekaké, $x$ ngwakili
\emph{kabèh} obyèk mawa sipat~$P$.

\begin{prop}
  Kanggo sembarang himpunan $A$ lan $B$, $A \subseteq A \cup B$.
\end{prop}

\begin{proof}
  Anggepen $A$ lan $B$ sembarang himpunan. Kita arep mbuktekaké
  $A \subseteq A \cup B$. Miturut dhéfinisi $\subseteq$, iki
  ateges: kanggo saben $x$, yèn $x \in A$ mula $x \in A \cup B$.
  Mula jupuken $x \in A$ minangka !!{element} sembarang saka~$A$.
  Kang kudu dibuktèkaké yaiku $x \in A \cup B$. Amarga $x \in A$,
  kita ngerti $x \in A$ utawa $x \in B$. Mula $x \in
  \Setabs{x}{x \in A \lor x \in B}$. Miturut dhéfinisi $\cup$,
  iki ateges $x \in A \cup B$.
\end{proof}

\subsection{Bukti Mawa Kasus}

Anggepen ana disjungsi minangka anggepan utawa kesimpulan kang
wis kabukten: $p$ utawa $q$ bener. Panjenengan arep mbuktekaké
$r$. Lakonana rong langkah: anggepen $p$ bener lan buktèkna~$r$;
banjur anggepen $q$ bener lan buktèkna~$r$ maneh. Cara iki sah
amarga salah siji saka rong kamungkinan mesthi kelakon. Rong
langkah mau nuduhaké yèn saben kamungkinan cukup kanggo
mbuktekaké~$r$. Yèn kaloroné bener, kita mung nduwèni rong
alasan yèn $r$ bener; ora perlu mbuktekaké $r$ maneh ing
anggepan yèn $p$ lan~$q$ padha bener. Kanggo nandhani cara iki,
kita ngumumaké yèn bakal ``mbédakaké kasus''. Upamané, anggepen
kita ngerti $x \in B \cup C$. Himpunan $B \cup C$
didhéfinisikaké minangka $\Setabs{x}{x \in B \text{ utawa } x \in C}$.
Mula miturut dhéfinisi, $x \in B$ utawa $x \in C$.
Kanggo mbuktekaké $x \in A$ saka iki, dhisik anggepen $x \in B$
lan buktèkna $x \in A$ saka anggepan iku; banjur anggepen
$x \in C$ lan buktèkna maneh $x \in A$. Ing ngisor anggepan
tulisen ``Kita mbédakaké kasus'', banjur ``Kasus (1): $x \in B$'',
lan ing tengah kertas ``Kasus (2): $x \in C$''. Sabanjuré
isinen pérangan ndhuwur lan ngisor kertas.

Bukti mawa kasus mligi migunani yèn kang dibuktèkaké dhéwé
awujud disjungsi. Iki conto prasaja:

\begin{prop}
Anggepen $B \subseteq D$ lan $C \subseteq E$.
Mula $B \cup C \subseteq D \cup E$.
\end{prop}

\begin{proof}
  Anggepen (a) $B \subseteq D$ lan (b) $C \subseteq E$.
  Miturut dhéfinisi, saben $x \in B$ uga $x \in D$ (c),
  lan saben $x \in C$ uga $x \in E$ (d).
  Kanggo mbuktekaké $B \cup C \subseteq D \cup E$, kita
  kudu nuduhaké yèn $x \in B \cup C$ mula $x \in D \cup E$
  (miturut dhéfinisi $\subseteq$). $x \in B \cup C$ yèn
  lan mung yèn $x \in B$ utawa $x \in C$ (miturut dhéfinisi~$\cup$).
  Semono uga $x \in D \cup E$ yèn lan mung yèn $x \in D$
  utawa $x \in E$. Mula kang kudu dibuktèkaké yaiku:
  kanggo sembarang $x$, yèn $x \in B$ utawa $x \in C$,
  mula $x \in D$ utawa $x \in E$.

  \begin{quote}
  Nganti kéné kita mung njlentrehaké dhéfinisi!{} Proposisi wis
  dipratelakaké maneh tanpa $\subseteq$ lan $\cup$, lan saiki
  kari bukti pratelan kondisional universal. Kaya kang wis dirembug,
  kita nindakaké iki kanthi nganggep sawijining $x$ kang nyukupi
  pérangan ``yèn'', banjur mbuktekaké pérangan ``mula''.
  Tegesé, kita nganggep $x \in B$ utawa $x \in C$ lan
  nuduhaké $x \in D$ utawa $x \in E$.\footnote{Alinea iki mung
  njlentrehaké apa kang lagi ditindakaké; dudu pérangan bukti,
  lan panjenengan ora perlu nulis kanthi rinci kaya iki nalika
  nyusun bukti dhéwé.}
  \end{quote}

  Anggepen $x \in B$ utawa $x \in C$. Kang kudu dibuktèkaké
  yaiku $x \in D$ utawa $x \in E$. Kita mbédakaké kasus.

  Kasus 1: $x \in B$. Miturut (c), $x \in D$. Mula $x \in D$
  utawa $x \in E$. Iki inferensi kang wis dirembug ing
  pérangan sadurungé.

  Kasus 2: $x \in C$. Miturut (d), $x \in E$. Mula $x \in D$
  utawa $x \in E$.
\end{proof}

\subsection{Mbuktekaké Pratelan Eksistensial}

Yèn dijaluk mbuktekaké pratelan eksistensial, pitakonané lumrahé
awujud ``buktèkna yèn ana~$x$ kang nyukupi $\dots x \dots$'',
yaiku ana obyèk mawa sipat kang diterangaké déning ``$\dots x
\dots$''. Ing kéné panjenengan kudu nemokaké obyèk kang cocog
lan nuduhaké yèn obyèk mau nduwèni sipat kang dijaluk.
Iki katon prasaja, nanging bukti kaya iki bisa ruwet. Asring
kita kudu \emph{ngonstruksi} utawa \emph{ndhéfinisikaké} obyèk
lan mbuktekaké yèn obyèk kang didhéfinisikaké iku nduwèni
sipat kang dikarepaké. Nemokaké obyèk kang pas bisa angel;
mbuktekaké sipaté uga bisa angel; malah kadhangkala nuduhaké
yèn dhéfinisiné pancèn ngasilaké obyèk uga angel!{}

Lumrahé panjenengan nemtokaké obyèk iku kanthi nulis,
upamané ``tetepna $x$ minangka \dots'' (ing kéné \dots{}
nerangaké obyèk kang dikarepaké), mbokmenawa mbuktekaké
yèn $\dots$ pancèn nemtokaké obyèk kang ana, banjur nuduhaké
yèn $x$ nduwèni sipat~$Q$. Iki conto prasaja.

\begin{prop}
  Anggepen $x \in B$. Mula ana~$A$ kang nyukupi
  $A \subseteq B$ lan $A \neq \emptyset$.
\end{prop}

\begin{proof}
  Anggepen $x \in B$. Tetepna $A = \{x\}$.
  \begin{quote}
    Ing kéné kita ndhéfinisikaké himpunan~$A$ kanthi nyebut
    !!{element}s siji-siji. Amarga wis dianggep yèn $x$ iku obyèk,
    lan himpunan bisa kawangun saka dhaptar !!{element}s-é,
    ora perlu mbuktekaké maneh yèn kita kasil ndhéfinisikaké
    himpunan~$A$. Nanging kita isih kudu nuduhaké yèn $A$
    nduwèni sipat-sipat kang dijaluk déning proposisi.
    Tanpa iku, buktiné durung jangkep!{}
  \end{quote}
  Amarga $x \in A$, $A \neq \emptyset$.
  \begin{quote}
    Iki gumantung marang dhéfinisi $A$ minangka $\{x\}$ lan
    kasunyatan cetha yèn $x \in \{x\}$ lan $x \notin \emptyset$.
  \end{quote}
  Amarga $x$ siji-sijiné !!{element} saka~$\{x\}$ lan
  $x \in B$, saben !!{element} saka~$A$ uga !!a{element}
  saka~$B$. Miturut dhéfinisi~$\subseteq$, $A \subseteq B$.
\end{proof}

\subsection{Migunakaké Pratelan Eksistensial}

Anggepen panjenengan ngerti sawijining pratelan eksistensial
bener, amarga wis dibuktèkaké utawa kalebu hipotesis kang bisa
digunakaké, upamané ``kanggo sawijining~$x$, $x \in A$'' utawa
``ana $x \in A$''. Kanggo migunakaké ing bukti, panjenengan
bisa ngenalaké jeneng kanggo salah siji obyèk kang ana miturut
hipotesis. Amarga $A$ paling ora ngemot siji obyèk, jeneng mau
bisa nuding salah siji obyèk iku. Bisa waé panjenengan ora bisa
milih utawa njlèntrèhaké obyèk mau luwih rinci, kajaba yèn
obyèk mau ana ing $A$. Nanging kanggo bukti, anggepen obyèk
iku wis dipilih lan wènèhana jeneng. Jeneng anyar aja nganti
wis dienggo sadurungé utawa muncul ing hipotesis, amarga bisa
nimbulaké panalaran kang klèru. Ing bukti, iki ditandhani
kanthi owah saka ``kanggo sawijining $x$, $x \in A$'' dadi
``Jupuken $a \in A$''. Saiki panjenengan bisa nalar babagan~$a$
lan migunakaké hipotesis liya nganti tekan kesimpulan $p$.
Yèn $p$ ora nyebut~$a$ maneh, $p$ ora gumantung
marang pilihan saksi $a \in A$ kang tartamtu,
saéngga katut saka anggepan
``kanggo sawijining $x$, $x \in A$''.

\begin{prop}
Yèn $A \neq \emptyset$, mula $A \cup B \neq \emptyset$.
\end{prop}

\begin{proof}
  Anggepen $A \neq \emptyset$. Dadi kanggo sawijining $x$,
  $x \in A$.
  \begin{quote}
    Ing kéné kita luwih dhisik mung mratelakaké hipotesis
    proposisi kanthi cara liya. Hipotesis $A \neq \emptyset$
    nyimpen pratelan eksistensial kang katon sawisé sawetara
    dhéfinisi dijlentrehaké. Dhéfinisi $=$ nuduhaké yèn
    $A = \emptyset$ yèn lan mung yèn saben $x \in A$
    uga $x \in \emptyset$ lan saben $x \in \emptyset$
    uga $x \in A$. Kanthi negasi, kita olèh yèn $A \neq
    \emptyset$ yèn salah siji ana $x \in A$ kang
    $x \notin \emptyset$, utawa ana $x \in \emptyset$
    kang $x \notin A$. Amarga ora ana $x \in \emptyset$,
    pérangan kapindho ora mungkin bener, déné ``$x \in A$
    lan $x \notin \emptyset$'' mung dadi $x \in A$.
    Dadi $A \neq \emptyset$ yèn lan mung yèn ana $x$
    kang $x \in A$. Iki pratelan eksistensial. Saiki
    kita migunakaké pratelan iku kanthi ngenalaké jeneng
    kanggo salah siji !!{element}s saka~$A$:
  \end{quote}
  Jupuken $a \in A$.
  \begin{quote}
    Saiki kita wis ngenalaké jeneng kanggo salah siji
    obyèk ing~$A$. Kita bakal nalar babagan~$a$, nanging
    mung nganggep yèn $a \in A$ lan ora luwih:
  \end{quote}
  Amarga $a \in A$, mula $a \in A \cup B$ miturut
  dhéfinisi~$\cup$. Dadi ana $x$ kang $x \in A \cup B$,
  yaiku $A \cup B \neq \emptyset$.
  \begin{quote}
    Ing langkah pungkasan, kita pindhah saka ``$a \in A \cup B$''
    menyang ``ana $x$ kang $x \in A \cup B$''.
    Pratelan pungkasan ora nyebut $a$ maneh, mula
    ``ana $x$ kang $x \in A \cup B$'' katut mung saka
    ``ana $x$ kang $x \in A$''. Iki ateges $A \cup B \neq
    \emptyset$.
  \end{quote}
\end{proof}

Mbokmenawa prayoga misahaké variabel kaiket kaya ``$x$''
saka jeneng saksi hipotètis kaya $a$, kaya ing ndhuwur.
Nanging ing laku kita asring migunakaké $x$ waé, kaya iki:
\begin{quote}
Anggepen $A \neq \emptyset$, yaiku ana $x \in A$.
Miturut dhéfinisi $\cup$, $x \in A \cup B$.
Dadi $A \cup B \neq \emptyset$.
\end{quote}
Yèn nganggo cara iki, panjenengan kudu luwih ngati-ati
supaya $x$ lan $y$ kang beda digunakaké kanggo pratelan
eksistensial kang béda. Upamané, iki \emph{dudu} bukti
kang sah kanggo ``Yèn $A \neq \emptyset$ lan $B \neq
\emptyset$, mula $A \cap B \neq \emptyset$''; pratelan
iku pancèn salah.
\begin{quote}
Anggepen $A \neq \emptyset$ lan $B \neq \emptyset$.
Dadi kanggo sawijining $x$, $x \in A$, lan kanggo
sawijining $x$, $x \in B$. Amarga $x \in A$ lan $x \in
B$, mula $x \in A \cap B$ miturut dhéfinisi~$\cap$.
Dadi $A \cap B \neq \emptyset$.
\end{quote}
Apa panjenengan bisa nemokaké langkah kang klèru lan nerangaké
ngapa asil mau ora bener?

\end{document}
